\PassOptionsToPackage{table,xcdraw}{xcolor}
\documentclass[11pt, a4paper]{article}

\usepackage[T1]{fontenc}
\usepackage[utf8]{inputenc}
\usepackage{tgpagella}
\usepackage[spanish, es-noshorthands]{babel}
\usepackage{amsmath, amssymb}
\usepackage{booktabs}
\usepackage{tabularx}
\usepackage[most]{tcolorbox}
\usepackage[margin=2.5cm]{geometry}
\usepackage{fancyhdr}

\pagestyle{fancy}
\fancyhf{}
\setlength{\headheight}{16pt}
\lhead{\textbf{UCB Sede Tarija} | Documento Docente}
\chead{Suplemento Técnico: Lab 3}
\rhead{Semana 8 - Sesión 2}
\lfoot{Circuitos Electrónicos II}
\rfoot{Página \thepage}
\renewcommand{\headrulewidth}{0.4pt}
\definecolor{ucbblue}{RGB}{0, 51, 102}

\begin{document}

\begin{center}
    \Large\textbf{\textcolor{ucbblue}{Suplemento del Instructor — Lab 3: Caracterización DC Ampliada}}[cite: 14]
\end{center}

\begin{tcolorbox}[colback=white, colframe=black, title=\textbf{Propósito de este Suplemento}]
    Complementar, \textbf{no reemplazar}, los artefactos de Lab 3 ya distribuidos[cite: 14]. Introduce dos modificaciones principales:
    \begin{enumerate}
        \item Alimentación nominal: $\boxed{\pm15\,\text{V} \rightarrow \pm10\,\text{V}}$[cite: 14].
        \item Prueba diferencial: $\boxed{\text{Un único punto} \rightarrow \text{Pequeño barrido DC}}$[cite: 14].
    \end{enumerate}
\end{tcolorbox}

\section*{1. Justificación Técnica de la Alimentación a $\pm10\,\text{V}$[cite: 14]}
Para el LM741 de Texas Instruments, el rango de tensión total de alimentación publicado comienza en $10\,\text{V}$ y el dispositivo no es \textit{rail-to-rail}[cite: 14]. El laboratorio utiliza salidas previstas de $0$ a $1\,\text{V}$[cite: 14]. Operar a $\pm10\,\text{V}$ (total de $20\,\text{V}$) proporciona un margen amplio sin mantener innecesariamente $30\,\text{V}$ en la protoboard[cite: 14]. No se recomienda reducir a $\pm5\,\text{V}$ ya que el margen respecto a los rieles sería demasiado limitante[cite: 14].

\section*{2. Lectura Pedagógica de los Resultados Experimentales[cite: 14]}
El nuevo objetivo es: \textit{Caracterizar experimentalmente la respuesta DC, identificar discrepancias y comprobar qué componente del error se reduce por compensación}[cite: 14].

\textbf{Análisis Gráfico y Analítico Esperado en Estudiantes[cite: 14]:}
\begin{itemize}
    \item \textbf{Línea recta aproximada con intercepto distinto de cero:} La red tiene la ganancia prevista, pero existe un offset. No asuma que todo el intercepto es $V_{OS}$ sin descartar mismatch resistivo o corrientes de polarización ($I_B$)[cite: 14].
    \item \textbf{Pendiente incorrecta (Error de ganancia):} Priorice revisar valores de resistores, ratios (mismatch), topología incorrecta o efectos de carga de la fuente antes de culpar al LM741[cite: 14].
    \item \textbf{Puntos fuertemente erráticos:} Problema de cableado, alimentación, referencias de tierra sueltas o mal uso del instrumento[cite: 14].
\end{itemize}

\section*{3. Uso y Efecto de la Compensación[cite: 14]}
Aplicar el ajuste offset-null (según datasheet) bajo $V_1=V_2=0\,\text{V}$[cite: 14]. \\
La observación pedagógicamente deseable \textbf{no es} que "todo quede perfecto en $0\,\text{V}$ en todo el rango", sino que el estudiante descubra lo siguiente[cite: 14]:
\begin{quote}
    \textit{Compensar el offset puede mover principalmente el intercepto (corrige el error DC interno), pero no necesariamente corrige una pendiente anómala (causada por tolerancias externas en la red de resistores).}[cite: 14]
\end{quote}
Esto ayuda a que el estudiante distinga empíricamente entre \texttt{internal device error} y \texttt{network ratio error} sin requerir todavía un análisis completo de presupuesto de errores (error budget)[cite: 14].

\section*{4. LTSpice y Evidencia Mínima Aceptable[cite: 14]}
\textbf{Prioridad Absoluta:} $\boxed{\text{Physical Evidence} > \text{Simulation}}$[cite: 14]. \\
LTSpice es una actividad suplementaria para contrastar $V_{o,\mathrm{ideal}}$ vs $V_{o,\mathrm{sim}}$ vs $V_{o,\mathrm{physical}}$[cite: 14]. Puede completarse de manera asincrónica si el tiempo no alcanza en la sesión física[cite: 14].

\textbf{Evidencia Mínima para un Laboratorio Significativo[cite: 14]:}
\begin{itemize}
    \item Circuito construido y verificado[cite: 14].
    \item Medición de al menos los puntos D0, D2 y D3[cite: 14].
    \item Discrepancia calculada e intento de troubleshooting documentado[cite: 14].
    \item Compensación de offset intentada (si el hardware lo permitió)[cite: 14].
    \item Una comparación Before/After documentada, junto con una interpretación final[cite: 14].
\end{itemize}

\section*{5. Gestión del Tiempo en Sesión (105 min)[cite: 14]}
\begin{tabularx}{\textwidth}{@{} l X @{}}
    \textbf{5–10 min:} & Readiness + prediction[cite: 14] \\
    \textbf{20 min:} & Construcción física + pre-power check[cite: 14] \\
    \textbf{10–15 min:} & Barrido D0–D3 (Baseline)[cite: 14] \\
    \textbf{15–20 min:} & Diagnóstico "Just-in-Time" ($V_{OS}, I_B, I_{OS}$, mismatch)[cite: 14] \\
    \textbf{10–15 min:} & Compensación de offset-null[cite: 14] \\
    \textbf{10 min:} & Repetición de puntos clave o modo común[cite: 14] \\
    \textbf{5–10 min:} & Interpretación y captura de evidencia[cite: 14]
\end{tabularx}

\end{document}
